\begin{frame}{역사: Sohl-Dickstein 2015 $\to$ DDPM 2020}
  \begin{itemize}
    \item 2015년 \kb{Sohl-Dickstein} 외 ``Deep Unsupervised Learning
      using Nonequilibrium Thermodynamics'' — 정방향/역방향 아이디어
      \textbf{처음 딥러닝에 제안}. 하지만 \textbf{학습이 너무 느려}
      실용화 못 함
    \item 2020년 \kb{DDPM} (Ho, Jain \& Abbeel, ``Denoising
      Diffusion Probabilistic Models''):
      (i) 노이즈 예측 매개변수화 $\epsilon_\theta$,
      (ii) 단순한 MSE 손실,
      (iii) 실용적 스케줄
      → \textbf{실제로 쓸 만한} 알고리즘으로 만듦
    \item 같은 해 \kb{DDIM} (Song et al.): 역방향 과정을
      \textbf{거의 결정적으로} 만들어 $T$를 \textbf{1000에서 약 50}
      단계로 줄임
    \item 2022년 \kb{Stable Diffusion} (Stability AI): 여기에 15.1절의
      \textbf{VAE}를 조합 — 이미지를 VAE 인코더로 \textbf{8$\times$
      압축한 잠재 공간}에서 Diffusion을 돌림, 계산량 $1/8^2$로
      감소
    \item 15.1에서 ``VAE의 매끄러운 잠재 공간''이 왜 중요한지 —
      그것이 곧 Diffusion을 \textbf{안정 있게} 학습·추론할 수 있게
      하는 \textbf{전제 조건}
    \item 같은 해 DALL$\cdot$E 2, Midjourney, 2024년 Sora(영상)까지
      — 이미지 생성의 \textbf{주류가 Diffusion으로 이동}
  \end{itemize}
  \vskip0.5em
  \kq{세 원리가 각자 독립적으로 끝나는 것이 아니라 \textbf{서로
  연결}되어 실전 시스템으로 합쳐진다는 점 — 이번 장을 ``세 개를
  나란히'' 보는 이유.}
\end{frame}
